\chapter{Anal Cancer}

\section*{Definition and Overview}
Anal cancer is a rare cancer that arises in the cells of the anus, the short canal at the end of the digestive tract through which stool leaves the body. Most anal cancers begin in the squamous cells that line the anal canal and are therefore called squamous cell carcinomas. Other, rarer types arise from glandular cells, skin pigment cells, or other tissues.

Anal cancer is distinct from the more common colorectal cancer, which develops higher up in the large bowel. When found early, anal cancer is often highly treatable, and many people are cured with a combination of chemotherapy and radiotherapy without needing surgery.

\section*{Epidemiology}
Anal cancer is uncommon, with roughly 1 in 500 to 1 in 1,000 people developing it in their lifetime. The incidence is higher in women than in men overall, although rates are considerably higher in men who have sex with men, particularly those living with HIV. Most cases occur in people over the age of 50.

The incidence has been rising in many developed countries over recent decades, largely linked to infection with human papillomavirus (HPV), which causes most cases. Anal cancer nevertheless remains far less common than colorectal cancer.

\section*{Aetiology and Pathophysiology}
Most anal cancers develop after long-standing infection with certain high-risk types of human papillomavirus (HPV), especially HPV type 16. The virus can cause pre-cancerous changes in the cells lining the anus that may, over many years, progress to invasive cancer. Factors that increase the risk include:
\begin{itemize}
    \item Persistent infection with high-risk HPV types
    \item A weakened immune system, such as from HIV, organ transplantation, or immunosuppressive medicines
    \item Anal receptive intercourse and multiple sexual partners
    \item Smoking, which increases the risk and worsens outcomes
    \item A history of other HPV-related cancers, such as cervical or vulval cancer
\end{itemize}
Once a cell becomes cancerous, it can grow into the surrounding tissue and spread to the lymph nodes in the groin and pelvis, and eventually to other organs.

\section*{Clinical Features}
The symptoms of anal cancer are often mistaken for common benign conditions such as haemorrhoids, so persistent symptoms should always be investigated:
\begin{itemize}
    \item Bleeding from the back passage, often fresh and bright red
    \item Pain, discomfort, or a feeling of pressure in the anal area
    \item A lump or mass at the anal opening
    \item Itching or a change in skin texture around the anus
    \item Changes in bowel habit, such as a feeling of not emptying the bowels
    \item Narrowing of the stools or increasing difficulty passing them
    \item Discharge of mucus or fluid from the anus
    \item A wound, ulcer, or fissure that does not heal
    \item Less commonly, lumps in the groin from spread to lymph nodes
\end{itemize}

\section*{Differential Diagnosis}
Many benign conditions produce similar symptoms and are far more common:
\begin{itemize}
    \item Haemorrhoids (piles)
    \item Anal fissures
    \item Anal fistulas and perianal abscesses
    \item Anal warts (genital warts) and pre-cancerous changes (anal intraepithelial neoplasia)
    \item Pruritus ani, or persistent anal itching
    \item Inflammatory bowel disease, such as Crohn's disease
    \item Other cancers in the region, including colorectal cancer or spread from cancers elsewhere
\end{itemize}

\section*{When to Seek Medical Care}
See a doctor for any anal symptom that does not settle within a few weeks, especially:
\begin{itemize}
    \item Bleeding from the back passage
    \item A persistent lump, mass, or growth at the anus
    \item Anal pain or discomfort that is getting worse
    \item A non-healing wound or ulcer
    \item Changes in bowel habit or narrowing of the stools
    \item Unexplained weight loss or fatigue
\end{itemize}

\textbf{Call 000 (or local emergency services) for:}
\begin{itemize}
    \item Heavy bleeding from the back passage
    \item Severe abdominal or anal pain
    \item A sudden inability to pass urine or stool
\end{itemize}

\section*{Diagnosis and Investigations}
\begin{itemize}
    \item \textbf{History and examination:} inspection of the anus and a digital rectal examination to feel for a mass, together with examination of the groins for lymph nodes
    \item \textbf{Proctoscopy or anoscopy:} direct viewing of the anal canal with a small instrument
    \item \textbf{Biopsy:} a small sample of the lesion is taken for laboratory examination, the only way to confirm cancer
    \item \textbf{Imaging:} MRI, CT, or endo-anal ultrasound to assess the size of the tumour and whether it has spread
    \item \textbf{HPV testing:} may be performed on biopsy samples
    \item \textbf{Blood tests and immunity checks:} as part of the overall assessment, including HIV testing where appropriate
\end{itemize}

\section*{Management}
Treatment is planned by a multidisciplinary team and depends on the stage, size, and type of the cancer:
\begin{itemize}
    \item \textbf{Chemoradiotherapy:} a combination of chemotherapy and radiotherapy is the standard treatment for most anal cancers and usually avoids surgery
    \item \textbf{Surgery:} removal of the tumour for small early cancers; for cancers that persist or return, removal of the anus and rectum (abdominoperineal resection) with a stoma may be needed
    \item \textbf{Surgery to the groin:} removal of affected lymph nodes in some cases
    \item \textbf{Supportive care:} management of pain, skin reactions, and bowel symptoms during and after treatment
    \item \textbf{Follow-up:} regular examinations and imaging for several years to detect recurrence early
\end{itemize}

\section*{Complications}
\begin{itemize}
    \item Side effects of treatment, including skin burns and soreness around the anus, diarrhoea, fatigue, and bladder irritation from radiotherapy
    \item Long-term bowel problems after radiotherapy, including urgency, leakage, and strictures (narrowing)
    \item Persistent or recurrent disease requiring more aggressive treatment
    \item Spread (metastasis) to the lymph nodes, liver, or lungs in advanced disease
    \item The need for a permanent stoma if the anus and rectum must be removed
    \item Sexual and psychological effects of the diagnosis and its treatment
\end{itemize}

\section*{Prognosis and Outcomes}
Anal cancer is one of the more treatable cancers, especially when it is confined to the anus. Most people with early anal cancer are cured, with five-year survival well over half overall and considerably higher for early disease. Chemoradiotherapy cures many people without the need for surgery, and a stoma is required only in a minority. Outcomes are less favourable if the cancer has spread to the lymph nodes or distant organs, or if the person has a weakened immune system. Regular follow-up is important because the cancer can return, sometimes years later.

\section*{Prevention and Self-Care}
\begin{itemize}
    \item The HPV vaccine protects against the virus types that cause most anal cancers; ask your doctor whether it is recommended for you
    \item Attend regular check-ups and screening if you are at higher risk, such as living with HIV or having had other HPV-related cancers
    \item Quit smoking, as it increases the risk of anal cancer and of recurrence
    \item Practise safer sex and consider condom use to reduce exposure to HPV
    \item Do not assume that anal bleeding is just haemorrhoids; report it to a doctor
    \item If you have anal warts, have them assessed and treated
    \item Maintain a healthy weight and general health to support the immune system
\end{itemize}

\section*{Sources and Further Reading}
The following resources provide reliable, peer-reviewed and patient-orientated information on anal cancer.

\begin{itemize}
    \item NHS. \textit{Overview of anal cancer, symptoms, and treatment.} https://www.nhs.uk/conditions/
    \item Mayo Clinic. \textit{Guide to the diagnosis and treatment of anal cancer.} https://www.mayoclinic.org/
    \item Centers for Disease Control and Prevention. \textit{Information on HPV and HPV-related cancers, including anal cancer.} https://www.cdc.gov/
    \item MSD Manual. \textit{Professional and patient information on anal cancer.} https://www.msdmanuals.com/
    \item MedlinePlus. \textit{Health information on anal cancer.} https://medlineplus.gov/
\end{itemize}
